\scanpage{13}
On peut aussi l'interpréter, \inserted{via $(P_i)_{\unsure{i\in\tau|i'\in\tau|i\in T}}\longmapsto\coprod_{i\in\tau}P_i$,} comme la catégorie des ensembles $E$ de cardinal $6$, munis d'une partition \inserted{$(P_i)_{i\in\tau}$} de type $(3,3)$ indexée par $\tau$, et pour chaque $i\in\tau$ d'une bijection $\omega_{P_i}\simeq D_i^*$. Ainsi les deux équivalences (7) fournissent une équivalence

\diagram{\begin{tikzcd}[ampersand replacement=\&,column sep=large]
{\llap{\text{(8)}\quad}\begin{array}[t]{@{}l@{}}
 \text{Catégorie des ens. }E\text{ de}\\
 \text{cardinal }6,\text{ munis d'une}\\
 \text{partition }\inserted{(P_i)_{i\in\tau}}\text{ de type }(3,3),\\
 \text{indexée par }\tau,\text{ et d'une}\\
 \text{\deleted{\unsure{f|famille} paire} de bijections}\\
 D_i^*\simeq\omega_{P_i}
\end{array}} \arrow[r,"\approx"] \&
{\begin{array}[t]{@{}l@{}}
 \text{Catégorie analogue}\\
 \text{définie en termes}\\
 \text{de }\tau'\text{ et de}\\
 (D_i^*)_{i\in\tau'}
\end{array}}
\end{tikzcd}}

Quand \inserted{(\unsure{gardant|grâce à} $V$)} on échange les rôles de $\tau$ et de $\tau'$, il est évident que le diagramme d'équivalences (7) reste le même à symétrie près, et que l'équivalence (8) est donc remplacée par \deleted{\unsure{l'}}\inserted{une} équivalence, \inserted{canoniquement} quasi-inverse de la précédente.

Les catégories source et but dans (8) sont définies en termes de la famille, \deleted{\unsure{de}}indexée par $P(V)=\tau\amalg\tau'$, d'ensembles à deux éléments $\omega_i=D_i^*\deleted{\unsure{\omega}}$ (dont la donnée équivaut à celle de la famille des droites \deleted{\unsure{D}}\unsure{Vect.|vect.} $D_i$ sur $\mathbb F_3$),
% archive sheet number: 12
